\documentclass[aspectratio=169,11pt]{beamer}
% Lecture 03 — Conditionals and if/when expressions
% Course: Introduction to Programming with Kotlin

% Shared Beamer preamble for Introduction to Programming with Kotlin
% Include with: % Shared Beamer preamble for Introduction to Programming with Kotlin
% Include with: % Shared Beamer preamble for Introduction to Programming with Kotlin
% Include with: \input{preamble}

\usetheme{Madrid}
\usecolortheme{default}

% Clean, professional palette (deep navy + muted teal accent)
\definecolor{LectureNavy}{HTML}{1B3A4B}
\definecolor{LectureAccent}{HTML}{2A9D8F}
\definecolor{LectureGray}{HTML}{4A5568}
\definecolor{CodeBg}{HTML}{F7FAFC}
\definecolor{AlertBg}{HTML}{FFF5F5}
\definecolor{TryBg}{HTML}{F0FFF4}
\definecolor{QuizBg}{HTML}{EBF8FF}
\definecolor{AnswerKeyBg}{HTML}{FFFFF0}
\definecolor{AnswerGold}{HTML}{B7791F}
\definecolor{KeywordBlue}{HTML}{2B6CB0}
\definecolor{StringGreen}{HTML}{276749}
\definecolor{CommentGray}{HTML}{718096}

\setbeamercolor{structure}{fg=LectureNavy}
\setbeamercolor{palette primary}{bg=LectureNavy,fg=white}
\setbeamercolor{palette secondary}{bg=LectureAccent,fg=white}
\setbeamercolor{palette tertiary}{bg=LectureGray,fg=white}
% Madrid puts title/subtitle in a coloured box (inherits palette primary bg).
% Title text must be light on that dark box — not navy-on-navy.
\setbeamercolor{title}{fg=white,bg=LectureNavy}
\setbeamercolor{subtitle}{fg=white,bg=LectureNavy}
\setbeamercolor{author}{fg=LectureNavy}
\setbeamercolor{date}{fg=LectureGray}
\setbeamercolor{institute}{fg=LectureGray}
\setbeamercolor{frametitle}{bg=LectureNavy,fg=white}
\setbeamercolor{block title}{bg=LectureNavy,fg=white}
\setbeamercolor{block body}{bg=CodeBg,fg=black}
\setbeamercolor{block title alerted}{bg=LectureAccent,fg=white}
\setbeamercolor{block body alerted}{bg=TryBg,fg=black}
\setbeamercolor{block title example}{bg=LectureGray,fg=white}
\setbeamercolor{block body example}{bg=AlertBg,fg=black}

\setbeamertemplate{navigation symbols}{}
\setbeamertemplate{itemize items}[circle]
\setbeamertemplate{enumerate items}[default]

\usepackage[T1]{fontenc}
\usepackage[utf8]{inputenc}
\usepackage{lmodern}
\usepackage{booktabs}
\usepackage{tabularx}
\usepackage{graphicx}
\usepackage{hyperref}
\usepackage{listings}
\usepackage{xcolor}
\usepackage{tikz}
\usetikzlibrary{positioning,arrows.meta}

% listings: Kotlin without shell-escape (no built-in Kotlin language)
\lstdefinelanguage{Kotlin}{
  morekeywords={
    abstract,actual,annotation,as,break,by,catch,class,companion,const,
    constructor,continue,crossinline,data,do,dynamic,else,enum,expect,
    external,false,field,file,final,finally,for,fun,get,if,import,in,
    infix,init,inline,inner,interface,internal,is,lateinit,noinline,null,
    object,open,operator,out,override,package,private,protected,public,
    reified,return,sealed,set,super,suspend,tailrec,this,throw,true,try,
    typealias,typeof,val,var,when,where,while
  },
  sensitive=true,
  morecomment=[l]{//},
  morecomment=[s]{/*}{*/},
  morestring=[b]",
  morestring=[b]',
  morestring=[s]{"""}{"""},
}

\lstdefinestyle{kotlinstyle}{
  language=Kotlin,
  basicstyle=\ttfamily\small,
  keywordstyle=\color{KeywordBlue}\bfseries,
  commentstyle=\color{CommentGray}\itshape,
  stringstyle=\color{StringGreen},
  numbers=left,
  numberstyle=\tiny\color{CommentGray},
  numbersep=6pt,
  frame=single,
  rulecolor=\color{LectureGray!40},
  backgroundcolor=\color{CodeBg},
  breaklines=true,
  showstringspaces=false,
  tabsize=2,
  captionpos=b,
  morekeywords={readln,readLine,println,print,listOf,mutableListOf,toInt,toDouble,toLong,toShort,toByte,toChar,toString,split,size},
}
\lstset{style=kotlinstyle}

\newcommand{\kotlincode}[1]{\texttt{\small #1}}
\newcommand{\trythis}[1]{%
  \begin{alertblock}{Try this}
    #1
  \end{alertblock}%
}
\newcommand{\commonmistake}[1]{%
  \begin{exampleblock}{Common mistake}
    #1
  \end{exampleblock}%
}
% Formative in-class checks: question slide, then a dedicated Answer slide
\newcommand{\quizpause}{%
  \par\vspace{0.5em}
  {\small\textit{Pause --- answer (clicker / raise hand / neighbour) before the next slide.}}%
}
\newcommand{\quizbox}[1]{%
  \setbeamercolor{block title}{bg=KeywordBlue,fg=white}%
  \setbeamercolor{block body}{bg=QuizBg,fg=black}%
  \begin{block}{Quiz}
    #1
  \end{block}%
  \setbeamercolor{block title}{bg=LectureNavy,fg=white}%
  \setbeamercolor{block body}{bg=CodeBg,fg=black}%
}
\newcommand{\answerbox}[1]{%
  \setbeamercolor{block title}{bg=AnswerGold,fg=white}%
  \setbeamercolor{block body}{bg=AnswerKeyBg,fg=black}%
  \begin{block}{Answer}
    #1
  \end{block}%
  \setbeamercolor{block title}{bg=LectureNavy,fg=white}%
  \setbeamercolor{block body}{bg=CodeBg,fg=black}%
}

% Empty author/date placeholders for the lecturer to fill
\author{}
\date{}


\usetheme{Madrid}
\usecolortheme{default}

% Clean, professional palette (deep navy + muted teal accent)
\definecolor{LectureNavy}{HTML}{1B3A4B}
\definecolor{LectureAccent}{HTML}{2A9D8F}
\definecolor{LectureGray}{HTML}{4A5568}
\definecolor{CodeBg}{HTML}{F7FAFC}
\definecolor{AlertBg}{HTML}{FFF5F5}
\definecolor{TryBg}{HTML}{F0FFF4}
\definecolor{QuizBg}{HTML}{EBF8FF}
\definecolor{AnswerKeyBg}{HTML}{FFFFF0}
\definecolor{AnswerGold}{HTML}{B7791F}
\definecolor{KeywordBlue}{HTML}{2B6CB0}
\definecolor{StringGreen}{HTML}{276749}
\definecolor{CommentGray}{HTML}{718096}

\setbeamercolor{structure}{fg=LectureNavy}
\setbeamercolor{palette primary}{bg=LectureNavy,fg=white}
\setbeamercolor{palette secondary}{bg=LectureAccent,fg=white}
\setbeamercolor{palette tertiary}{bg=LectureGray,fg=white}
% Madrid puts title/subtitle in a coloured box (inherits palette primary bg).
% Title text must be light on that dark box — not navy-on-navy.
\setbeamercolor{title}{fg=white,bg=LectureNavy}
\setbeamercolor{subtitle}{fg=white,bg=LectureNavy}
\setbeamercolor{author}{fg=LectureNavy}
\setbeamercolor{date}{fg=LectureGray}
\setbeamercolor{institute}{fg=LectureGray}
\setbeamercolor{frametitle}{bg=LectureNavy,fg=white}
\setbeamercolor{block title}{bg=LectureNavy,fg=white}
\setbeamercolor{block body}{bg=CodeBg,fg=black}
\setbeamercolor{block title alerted}{bg=LectureAccent,fg=white}
\setbeamercolor{block body alerted}{bg=TryBg,fg=black}
\setbeamercolor{block title example}{bg=LectureGray,fg=white}
\setbeamercolor{block body example}{bg=AlertBg,fg=black}

\setbeamertemplate{navigation symbols}{}
\setbeamertemplate{itemize items}[circle]
\setbeamertemplate{enumerate items}[default]

\usepackage[T1]{fontenc}
\usepackage[utf8]{inputenc}
\usepackage{lmodern}
\usepackage{booktabs}
\usepackage{tabularx}
\usepackage{graphicx}
\usepackage{hyperref}
\usepackage{listings}
\usepackage{xcolor}
\usepackage{tikz}
\usetikzlibrary{positioning,arrows.meta}

% listings: Kotlin without shell-escape (no built-in Kotlin language)
\lstdefinelanguage{Kotlin}{
  morekeywords={
    abstract,actual,annotation,as,break,by,catch,class,companion,const,
    constructor,continue,crossinline,data,do,dynamic,else,enum,expect,
    external,false,field,file,final,finally,for,fun,get,if,import,in,
    infix,init,inline,inner,interface,internal,is,lateinit,noinline,null,
    object,open,operator,out,override,package,private,protected,public,
    reified,return,sealed,set,super,suspend,tailrec,this,throw,true,try,
    typealias,typeof,val,var,when,where,while
  },
  sensitive=true,
  morecomment=[l]{//},
  morecomment=[s]{/*}{*/},
  morestring=[b]",
  morestring=[b]',
  morestring=[s]{"""}{"""},
}

\lstdefinestyle{kotlinstyle}{
  language=Kotlin,
  basicstyle=\ttfamily\small,
  keywordstyle=\color{KeywordBlue}\bfseries,
  commentstyle=\color{CommentGray}\itshape,
  stringstyle=\color{StringGreen},
  numbers=left,
  numberstyle=\tiny\color{CommentGray},
  numbersep=6pt,
  frame=single,
  rulecolor=\color{LectureGray!40},
  backgroundcolor=\color{CodeBg},
  breaklines=true,
  showstringspaces=false,
  tabsize=2,
  captionpos=b,
  morekeywords={readln,readLine,println,print,listOf,mutableListOf,toInt,toDouble,toLong,toShort,toByte,toChar,toString,split,size},
}
\lstset{style=kotlinstyle}

\newcommand{\kotlincode}[1]{\texttt{\small #1}}
\newcommand{\trythis}[1]{%
  \begin{alertblock}{Try this}
    #1
  \end{alertblock}%
}
\newcommand{\commonmistake}[1]{%
  \begin{exampleblock}{Common mistake}
    #1
  \end{exampleblock}%
}
% Formative in-class checks: question slide, then a dedicated Answer slide
\newcommand{\quizpause}{%
  \par\vspace{0.5em}
  {\small\textit{Pause --- answer (clicker / raise hand / neighbour) before the next slide.}}%
}
\newcommand{\quizbox}[1]{%
  \setbeamercolor{block title}{bg=KeywordBlue,fg=white}%
  \setbeamercolor{block body}{bg=QuizBg,fg=black}%
  \begin{block}{Quiz}
    #1
  \end{block}%
  \setbeamercolor{block title}{bg=LectureNavy,fg=white}%
  \setbeamercolor{block body}{bg=CodeBg,fg=black}%
}
\newcommand{\answerbox}[1]{%
  \setbeamercolor{block title}{bg=AnswerGold,fg=white}%
  \setbeamercolor{block body}{bg=AnswerKeyBg,fg=black}%
  \begin{block}{Answer}
    #1
  \end{block}%
  \setbeamercolor{block title}{bg=LectureNavy,fg=white}%
  \setbeamercolor{block body}{bg=CodeBg,fg=black}%
}

% Empty author/date placeholders for the lecturer to fill
\author{}
\date{}


\usetheme{Madrid}
\usecolortheme{default}

% Clean, professional palette (deep navy + muted teal accent)
\definecolor{LectureNavy}{HTML}{1B3A4B}
\definecolor{LectureAccent}{HTML}{2A9D8F}
\definecolor{LectureGray}{HTML}{4A5568}
\definecolor{CodeBg}{HTML}{F7FAFC}
\definecolor{AlertBg}{HTML}{FFF5F5}
\definecolor{TryBg}{HTML}{F0FFF4}
\definecolor{QuizBg}{HTML}{EBF8FF}
\definecolor{AnswerKeyBg}{HTML}{FFFFF0}
\definecolor{AnswerGold}{HTML}{B7791F}
\definecolor{KeywordBlue}{HTML}{2B6CB0}
\definecolor{StringGreen}{HTML}{276749}
\definecolor{CommentGray}{HTML}{718096}

\setbeamercolor{structure}{fg=LectureNavy}
\setbeamercolor{palette primary}{bg=LectureNavy,fg=white}
\setbeamercolor{palette secondary}{bg=LectureAccent,fg=white}
\setbeamercolor{palette tertiary}{bg=LectureGray,fg=white}
% Madrid puts title/subtitle in a coloured box (inherits palette primary bg).
% Title text must be light on that dark box — not navy-on-navy.
\setbeamercolor{title}{fg=white,bg=LectureNavy}
\setbeamercolor{subtitle}{fg=white,bg=LectureNavy}
\setbeamercolor{author}{fg=LectureNavy}
\setbeamercolor{date}{fg=LectureGray}
\setbeamercolor{institute}{fg=LectureGray}
\setbeamercolor{frametitle}{bg=LectureNavy,fg=white}
\setbeamercolor{block title}{bg=LectureNavy,fg=white}
\setbeamercolor{block body}{bg=CodeBg,fg=black}
\setbeamercolor{block title alerted}{bg=LectureAccent,fg=white}
\setbeamercolor{block body alerted}{bg=TryBg,fg=black}
\setbeamercolor{block title example}{bg=LectureGray,fg=white}
\setbeamercolor{block body example}{bg=AlertBg,fg=black}

\setbeamertemplate{navigation symbols}{}
\setbeamertemplate{itemize items}[circle]
\setbeamertemplate{enumerate items}[default]

\usepackage[T1]{fontenc}
\usepackage[utf8]{inputenc}
\usepackage{lmodern}
\usepackage{booktabs}
\usepackage{tabularx}
\usepackage{graphicx}
\usepackage{hyperref}
\usepackage{listings}
\usepackage{xcolor}
\usepackage{tikz}
\usetikzlibrary{positioning,arrows.meta}

% listings: Kotlin without shell-escape (no built-in Kotlin language)
\lstdefinelanguage{Kotlin}{
  morekeywords={
    abstract,actual,annotation,as,break,by,catch,class,companion,const,
    constructor,continue,crossinline,data,do,dynamic,else,enum,expect,
    external,false,field,file,final,finally,for,fun,get,if,import,in,
    infix,init,inline,inner,interface,internal,is,lateinit,noinline,null,
    object,open,operator,out,override,package,private,protected,public,
    reified,return,sealed,set,super,suspend,tailrec,this,throw,true,try,
    typealias,typeof,val,var,when,where,while
  },
  sensitive=true,
  morecomment=[l]{//},
  morecomment=[s]{/*}{*/},
  morestring=[b]",
  morestring=[b]',
  morestring=[s]{"""}{"""},
}

\lstdefinestyle{kotlinstyle}{
  language=Kotlin,
  basicstyle=\ttfamily\small,
  keywordstyle=\color{KeywordBlue}\bfseries,
  commentstyle=\color{CommentGray}\itshape,
  stringstyle=\color{StringGreen},
  numbers=left,
  numberstyle=\tiny\color{CommentGray},
  numbersep=6pt,
  frame=single,
  rulecolor=\color{LectureGray!40},
  backgroundcolor=\color{CodeBg},
  breaklines=true,
  showstringspaces=false,
  tabsize=2,
  captionpos=b,
  morekeywords={readln,readLine,println,print,listOf,mutableListOf,toInt,toDouble,toLong,toShort,toByte,toChar,toString,split,size},
}
\lstset{style=kotlinstyle}

\newcommand{\kotlincode}[1]{\texttt{\small #1}}
\newcommand{\trythis}[1]{%
  \begin{alertblock}{Try this}
    #1
  \end{alertblock}%
}
\newcommand{\commonmistake}[1]{%
  \begin{exampleblock}{Common mistake}
    #1
  \end{exampleblock}%
}
% Formative in-class checks: question slide, then a dedicated Answer slide
\newcommand{\quizpause}{%
  \par\vspace{0.5em}
  {\small\textit{Pause --- answer (clicker / raise hand / neighbour) before the next slide.}}%
}
\newcommand{\quizbox}[1]{%
  \setbeamercolor{block title}{bg=KeywordBlue,fg=white}%
  \setbeamercolor{block body}{bg=QuizBg,fg=black}%
  \begin{block}{Quiz}
    #1
  \end{block}%
  \setbeamercolor{block title}{bg=LectureNavy,fg=white}%
  \setbeamercolor{block body}{bg=CodeBg,fg=black}%
}
\newcommand{\answerbox}[1]{%
  \setbeamercolor{block title}{bg=AnswerGold,fg=white}%
  \setbeamercolor{block body}{bg=AnswerKeyBg,fg=black}%
  \begin{block}{Answer}
    #1
  \end{block}%
  \setbeamercolor{block title}{bg=LectureNavy,fg=white}%
  \setbeamercolor{block body}{bg=CodeBg,fg=black}%
}

% Empty author/date placeholders for the lecturer to fill
\author{}
\date{}


\title{Lecture 03: Conditionals and Expressions}
\subtitle{Introduction to Programming with Kotlin}

\begin{document}

%----------------------------------------------------------------------
    \begin{frame}
        \titlepage
    \end{frame}

%----------------------------------------------------------------------
    \begin{frame}{Agenda}
        \begin{enumerate}
            \item Conditional operator: \kotlincode{if} / \kotlincode{else}
            \item Combining conditions (Part 2)
            \item \kotlincode{if} and \kotlincode{when} as expressions
        \end{enumerate}
    \end{frame}

%======================================================================
    \section{Conditional operator}
%======================================================================

    \begin{frame}{Making decisions}
        Programs often choose different actions based on data.
        \vspace{0.6em}
        \begin{itemize}
            \item A \textbf{condition} is an expression that is \kotlincode{true} or \kotlincode{false}
              (type \kotlincode{Boolean})
            \item \kotlincode{if} runs a block when the condition is true
            \item \kotlincode{else} covers the other case
        \end{itemize}
    \end{frame}

    \begin{frame}{Comparison operators}
        \begin{tabular}{ll}
            \toprule
            Operator & Meaning \\
            \midrule
            \kotlincode{==} & equal to \\
            \kotlincode{!=} & not equal to \\
            \kotlincode{<} & less than \\
            \kotlincode{<=} & less than or equal \\
            \kotlincode{>} & greater than \\
            \kotlincode{>=} & greater than or equal \\
            \bottomrule
        \end{tabular}
        \vspace{0.6em}
        \commonmistake{Do not confuse \kotlincode{=} (assignment) with \kotlincode{==} (equality).}
    \end{frame}

    \begin{frame}[fragile]{Basic \kotlincode{if} / \kotlincode{else}}
        \begin{lstlisting}
fun main() {
    print("Enter score: ")
    val score = readln().toInt()
    if (score >= 50) {
        println("Pass")
    } else {
        println("Fail")
    }
}
        \end{lstlisting}
        \vspace{0.3em}
        Braces \kotlincode{\{ \}} mark the blocks. Prefer always using braces,
        even for one-liners.
    \end{frame}

    \begin{frame}[fragile]{\kotlincode{else if} chains}
        \begin{lstlisting}
fun main() {
    val temp = readln().toInt()
    if (temp < 0) {
        println("Freezing")
    } else if (temp < 15) {
        println("Cold")
    } else if (temp < 25) {
        println("Mild")
    } else {
        println("Warm")
    }
}
        \end{lstlisting}
        \vspace{0.3em}
        Conditions are checked top to bottom; the first true branch wins.
    \end{frame}

    \begin{frame}[fragile]{Nested \kotlincode{if} (use sparingly)}
        \begin{lstlisting}
fun main() {
    val age = readln().toInt()
    val hasTicket = readln() == "yes"
    if (age >= 18) {
        if (hasTicket) {
            println("Enter")
        } else {
            println("Buy a ticket")
        }
    } else {
        println("Too young")
    }
}
        \end{lstlisting}
        Nesting deep trees is hard to read --- combining conditions often helps.
    \end{frame}

    %--- Quiz checkpoint 1 ---
    \begin{frame}{Quiz 1 --- if / else}
        \quizbox{For \kotlincode{val x = 10}, which branch runs?\\
        \kotlincode{if (x > 10) A else if (x >= 10) B else C}}
        \begin{enumerate}
            \item A
            \item B
            \item C
            \item A and B
        \end{enumerate}
        \quizpause
    \end{frame}

    \begin{frame}{Answer --- Quiz 1}
        \answerbox{\textbf{2.} B}
        \vspace{0.4em}
        \kotlincode{x > 10} is false; \kotlincode{x >= 10} is true, so branch B runs.
        Only one branch of an \kotlincode{if}/\kotlincode{else if} chain executes.
    \end{frame}

%======================================================================
    \section{Combining conditions}
%======================================================================

    \begin{frame}{Logical operators}
        Combine Booleans with:
        \vspace{0.5em}
        \begin{tabular}{lll}
            \toprule
            Operator & Name & Idea \\
            \midrule
            \kotlincode{\&\&} & and & both must be true \\
            \kotlincode{||} & or & at least one true \\
            \kotlincode{!} & not & flips true/false \\
            \bottomrule
        \end{tabular}
        \vspace{0.6em}
        Kotlin uses short-circuit evaluation: if the result is already known,
        the rest is not evaluated.
    \end{frame}

    \begin{frame}[fragile]{And, or, not}
        \begin{lstlisting}
fun main() {
    val age = readln().toInt()
    val hasId = readln() == "yes"

    if (age >= 18 && hasId) {
        println("Adult with ID")
    }
    if (age < 13 || age > 65) {
        println("Discount age")
    }
    if (!(age >= 18)) {
        println("Under 18")
    }
}
        \end{lstlisting}
    \end{frame}

    \begin{frame}{Truth tables (quick)}
        \begin{columns}[T]
            \begin{column}{0.48\textwidth}
                \textbf{\kotlincode{\&\&} (and)}
                \begin{itemize}
                    \item T \&\& T $\rightarrow$ T
                    \item T \&\& F $\rightarrow$ F
                    \item F \&\& T $\rightarrow$ F
                    \item F \&\& F $\rightarrow$ F
                \end{itemize}
            \end{column}
            \begin{column}{0.48\textwidth}
                \textbf{\kotlincode{||} (or)}
                \begin{itemize}
                    \item T || T $\rightarrow$ T
                    \item T || F $\rightarrow$ T
                    \item F || T $\rightarrow$ T
                    \item F || F $\rightarrow$ F
                \end{itemize}
            \end{column}
        \end{columns}
        \vspace{0.6em}
        \kotlincode{!true} is \kotlincode{false}; \kotlincode{!false} is \kotlincode{true}.
    \end{frame}

    \begin{frame}[fragile]{Readable ranges}
        Check that a value lies in an interval:
        \vspace{0.3em}
        \begin{lstlisting}
fun main() {
    val n = readln().toInt()
    // Inclusive range check with &&
    if (n >= 1 && n <= 10) {
        println("Between 1 and 10")
    }
    // Kotlin also has `in` (more later with ranges)
    if (n in 1..10) {
        println("Also between 1 and 10")
    }
}
        \end{lstlisting}
        \trythis{Validate that a password length is at least 8 \emph{and} the user
        typed \kotlincode{yes} to confirm.}
    \end{frame}

    \begin{frame}{Operator precedence (practical)}
        \begin{itemize}
            \item Comparisons bind tighter than \kotlincode{\&\&} / \kotlincode{||}
            \item \kotlincode{\&\&} binds tighter than \kotlincode{||}
            \item When unsure, add parentheses
        \end{itemize}
        \vspace{0.5em}
        \begin{block}{Examples}
            \kotlincode{a > 0 \&\& b > 0 || c > 0}\\
            means \kotlincode{(a > 0 \&\& b > 0) || c > 0}\\[0.4em]
            Prefer: \kotlincode{(a > 0 \&\& b > 0) || (c > 0)}
        \end{block}
    \end{frame}

    %--- Quiz checkpoint 2 ---
    \begin{frame}{Quiz 2 --- Combining}
        \quizbox{When is \kotlincode{age >= 18 \&\& hasTicket} true?}
        \begin{enumerate}
            \item Only when age is at least 18
            \item Only when \kotlincode{hasTicket} is true
            \item When both are true
            \item When either is true
        \end{enumerate}
        \quizpause
    \end{frame}

    \begin{frame}{Answer --- Quiz 2}
        \answerbox{\textbf{3.} When both are true}
        \vspace{0.4em}
        \kotlincode{\&\&} requires both sides to be true.
        Use \kotlincode{||} if either condition is enough.
    \end{frame}

%======================================================================
    \section{if and when as expressions}
%======================================================================

    \begin{frame}{Statements vs expressions}
        In Kotlin, many control structures are \textbf{expressions}: they produce a value.
        \vspace{0.6em}
        \begin{itemize}
            \item You can assign the result of \kotlincode{if} to a variable
            \item Same idea with \kotlincode{when}
            \item This reduces duplicated \kotlincode{println} / assignment code
        \end{itemize}
    \end{frame}

    \begin{frame}[fragile]{\kotlincode{if} as an expression}
        \begin{lstlisting}
fun main() {
    val score = readln().toInt()
    val result = if (score >= 50) "Pass" else "Fail"
    println(result)

    val max = if (score > 100) 100 else score
    println("Capped: $max")
}
        \end{lstlisting}
        \vspace{0.3em}
        When used as an expression, \kotlincode{else} is required
        (both branches must produce a value).
    \end{frame}

    \begin{frame}[fragile]{Blocks as expression bodies}
        \begin{lstlisting}
fun main() {
    val x = readln().toInt()
    val label = if (x % 2 == 0) {
        println("Even path")
        "even"
    } else {
        println("Odd path")
        "odd"
    }
    println(label)
}
        \end{lstlisting}
        The last expression in the block is the value of that branch.
    \end{frame}

    \begin{frame}[fragile]{\kotlincode{when} as a better multi-branch}
        \begin{lstlisting}
fun main() {
    val code = readln().toInt()
    val mood = when (code) {
        1 -> "Happy"
        2 -> "Okay"
        3, 4 -> "Tired"
        else -> "Unknown"
    }
    println(mood)
}
        \end{lstlisting}
        \vspace{0.3em}
        \kotlincode{when} replaces long \kotlincode{else if} chains for discrete cases.
    \end{frame}

    \begin{frame}[fragile]{\kotlincode{when} with conditions}
        \begin{lstlisting}
fun main() {
    val n = readln().toInt()
    val kind = when {
        n < 0 -> "negative"
        n == 0 -> "zero"
        n % 2 == 0 -> "positive even"
        else -> "positive odd"
    }
    println(kind)
}
        \end{lstlisting}
        Subject-less \kotlincode{when \{ \}} checks Boolean conditions in order.
    \end{frame}

    \begin{frame}[fragile]{\kotlincode{when} with ranges (preview)}
        \begin{lstlisting}
fun main() {
    val score = readln().toInt()
    val grade = when (score) {
        in 90..100 -> "A"
        in 80..89 -> "B"
        in 70..79 -> "C"
        in 60..69 -> "D"
        else -> "F"
    }
    println(grade)
}
        \end{lstlisting}
        Ranges are covered in depth in Lecture~04.
    \end{frame}

    %--- Quiz checkpoint 3 ---
    \begin{frame}{Quiz 3 --- Expressions}
        \quizbox{Which is valid Kotlin?}
        \begin{enumerate}
            \item \kotlincode{val s = if (x > 0) "pos"}
            \item \kotlincode{val s = if (x > 0) "pos" else "nonpos"}
            \item \kotlincode{val s = when (x) \{ 1 -> "one" \}}
              without \kotlincode{else}, for arbitrary \kotlincode{Int x}
            \item \kotlincode{val s = if = (x > 0) "pos" else "no"}
        \end{enumerate}
        \quizpause
    \end{frame}

    \begin{frame}{Answer --- Quiz 3}
        \answerbox{\textbf{2.} \kotlincode{val s = if (x > 0) "pos" else "nonpos"}}
        \vspace{0.4em}
        An \kotlincode{if} expression needs \kotlincode{else}.
        A \kotlincode{when} expression over an \kotlincode{Int} needs to be exhaustive
        (usually via \kotlincode{else}). Option 4 is invalid syntax.
    \end{frame}

%======================================================================
    \section{Wrap-up}
%======================================================================

    \begin{frame}{Summary}
        \begin{itemize}
            \item \kotlincode{if} / \kotlincode{else if} / \kotlincode{else} choose branches
            \item Combine with \kotlincode{\&\&}, \kotlincode{||}, \kotlincode{!}
            \item Prefer parentheses for clarity
            \item \kotlincode{if} and \kotlincode{when} can produce values (expressions)
            \item Use \kotlincode{when} for multi-way matching
        \end{itemize}
    \end{frame}

    \begin{frame}{Looking ahead}
        Next: collections, \kotlincode{repeat} / \kotlincode{while}, \kotlincode{for} and
        \kotlincode{size}, ranges, and \kotlincode{do-while} with \kotlincode{split}.
    \end{frame}

\end{document}
